% ============================================================
% CONTRARIAN POSITIONS (insert between Design Principles and Limitations)
% ============================================================
\section{Engaging Contrarian Positions}
\label{sec:contrarian}

The framework's formal results invite strong objections from multiple intellectual
traditions.  Rather than constructing strawmen, we steelman seven contrarian
positions, group them thematically, and evaluate what the evidence actually
supports.  Each position draws on material originally developed across the
seven pillars; here we consolidate by theme rather than by pillar.  Table~\ref{tab:contrarian-summary} provides a summary; the subsections below
develop each position in detail.

\begin{table}[t]
\centering
\small
\begin{tabular}{lp{4.5cm}p{4.5cm}}
\toprule
\textbf{Position} & \textbf{Strongest Form} & \textbf{Framework Response} \\
\midrule
Scale makes governance obsolete & Scaling laws predict learned priors will outperform engineered methods & Abstraction gap is mechanism-agnostic; compound errors are mathematical identities \\
Formal methods are overkill & Situated action and language games challenge fixed specifications & $\sigma$ is a design parameter; formal methods at high $\sigma$, NL at low $\sigma$ \\
Governance is overhead & Cost exceeds benefit below a team-size threshold & Bifurcation theorem identifies exact threshold; AI velocity shifts it lower \\
ADRs are theater & No empirical evidence ADRs improve quality outcomes & Survival analysis with evidence expiry addresses staleness; empirical gap acknowledged \\
Humans should review everything & Human review is the gold standard & Throughput mismatch: 50\%+ of code is AI-generated; humans triage, not review exhaustively \\
AI quality is good enough & GitHub RCT shows no quality difference & Longitudinal data contradicts; accretion effects invisible to short-term studies \\
Theory cannot be preserved & Tacit knowledge has divergent description length & Lossy compression beats no compression; uncertainty markers address false confidence \\
\bottomrule
\end{tabular}
\caption{Summary of contrarian positions engaged, with their strongest forms and framework responses.}
\label{tab:contrarian-summary}
\end{table}

\subsection{The Bitter Lesson: Scale Makes Governance Obsolete}
\label{sec:contrarian-bitter}

\paragraph{The steelmanned argument.}
Sutton's Bitter Lesson~\citep{sutton2019bitter} observes that general methods
leveraging computation (search, learning) have historically defeated
hand-engineered, knowledge-intensive approaches.  LLMs trained on all publicly
available code are the epitome of the scaling approach.  If scaling continues to
improve, formal methods, governance mechanisms, and structured verification may
become unnecessary overhead, much as hand-tuned chess evaluation functions were
rendered obsolete by AlphaZero.  Dijkstra's programme of designing formal
languages for formal thought is precisely the kind of ``human-knowledge''
approach that scaling predicts will lose.

Distributed cognition theory~\citep{hutchins1995cognition, clark2003natural}
strengthens this position: the human + LLM + IDE + test suite constitutes a
distributed cognitive system that achieves precision even when no single
component (including the human's natural-language specification) is precise
alone.  Formally, the abstraction gap $\mathcal{G}(S_{\text{NL}}, P)$ may be
large, but $\mathcal{G}(S_{\text{NL}} \parallel \Phi_{\text{LLM}} \parallel
\mathcal{T}_{\text{tests}}, P)$ can be small under distributed composition.
This is a genuine strength of the natural-language approach that the framework
can accommodate but does not predict.

\paragraph{What the evidence shows.}
The framework is agnostic on mechanism.  If scaling obsoletes formal methods,
the abstraction gap $\mathcal{G}(S, P) = K(P \mid S)$ still exists; what
changes is the mechanism for traversing it (brute-force learned priors rather
than structured refinement).  The compound error sensitivity theorem
(Theorem~\ref{thm:sensitivity}) is a mathematical identity, not a claim about
any particular verification method: $R = \prod p_i$ holds whether the $p_i$ are
achieved by formal proofs, learned verifiers, or brute-force sampling.  The
Bitter Lesson predicts that \emph{how} we achieve high $p_i$ will shift; it
does not predict that the multiplicative structure of pipeline reliability will
change.

Moreover, current empirical evidence is mixed.  HumanEval pass rates exceed
90\%, but SWE-bench Verified plateaus around 77\%, with METR finding that
approximately 50\% of test-passing PRs would not be merged by human
reviewers~\citep{metr2026swebench}.  LiveCodeBench remains at 38.9\%.  The gap
between isolated-function benchmarks and real-world integration tasks has not
closed with scale alone.  The 50-point gap between HumanEval ($> 90\%$, low-NCD
tasks with short specifications) and LiveCodeBench (38.9\%, high-NCD tasks
requiring complex integration) is precisely what the abstraction gap framework
predicts: as task complexity grows, the LLM's learned prior can no longer
compensate for the information deficit in the specification.

The Goodhart taxonomy~\citep{manheim2019goodhart} provides additional grounds
for skepticism about pure-scaling approaches.  All four Goodhart mechanisms
(regressional, extremal, causal, adversarial) manifest in agent benchmarks:
\begin{itemize}
  \item \emph{Regressional:} optimising a noisy proxy (test pass rates) selects
    for noise in the proxy-target relationship.
  \item \emph{Extremal:} the proxy-target relationship, estimated in a typical
    regime, breaks down at extremes of task complexity.
  \item \emph{Causal:} intervening on the proxy (writing code to pass tests) can
    break the causal pathway to genuine quality.
  \item \emph{Adversarial:} agents can directly manipulate the proxy without
    affecting the target.
\end{itemize}
Every documented NIST CAISI cheating instance is detectable through structural
invariant violation checks, not through statistical proxy-target monitoring,
supporting the principle that structural enforcement addresses what scaling
alone cannot.  The ``true quality'' $Q$ in any gaming detection framework is
fundamentally uncomputable; all practical proxies inherit the Goodhart problem
they aim to detect.


\subsection{Formal Methods Are Overkill}
\label{sec:contrarian-formal}

\paragraph{The steelmanned argument.}
Suchman's situated action theory~\citep{suchman1987plans} argues that plans are
resources for action, not determinants of action.  If this is correct, the
refinement lattice (which models specification as a plan progressively
concretised into code) may be the wrong metaphor for iterative, conversational
development.  The specification-code relationship is not a static refinement
chain but a dynamic, context-dependent process where meaning emerges from
interaction.

Wittgenstein's later philosophy~\citep{wittgenstein1953philosophical} deepens
this critique.  If meaning is constituted by use rather than by formal
structure, then a prompt's informational content is not a fixed quantity
measurable by Kolmogorov complexity; it is a move in a language game whose
significance depends on shared practices.  This challenges a presupposition of
the entire information-theoretic framework: that specifications have fixed
informational content independent of context.

The democratisation argument adds practical weight: 5~billion natural-language
speakers versus 27~million programmers (a 185:1 ratio).  If natural-language
programming enables even 10\% of this population to direct computation, the
total value created may dwarf the quality loss on individual
programmes~\citep{litt2023malleable, altman2024intelligence}.  Prompt-based
rules achieve 77\% compliance (AgentSpec data), which may be ``good enough''
for most use cases.

\paragraph{What the evidence shows.}
These arguments are strongest when the specificity requirement is low
($\sigma < 0.3$) and weakest when formal guarantees are needed
($\sigma > 0.7$).  The framework accommodates both by treating $\sigma$ as a
design parameter, not a universal constant.  For prototypes and low-stakes
tooling, Layer~1 structural gates alone may suffice; for production systems with
safety constraints, the formal machinery becomes essential.

The 82\% success rate of vericoding~\citep{vericoding2025} (human writes formal
specification, AI generates implementation, formal verifier checks conformance)
demonstrates that formal methods can be made accessible through AI translation:
the human reviews a specification rather than authoring one from scratch.  This
addresses the historical cost barrier~\citep{kleppmann2025formal} without
abandoning formal guarantees.  The key insight is that AI can serve as a
\emph{formalism translator}, shifting the human role from specification author
to specification reviewer.  Reviewing a formal specification is vastly easier
than authoring one from scratch, which changes the cost calculus that
historically made formal methods prohibitive.

We acknowledge, following Suchman and Wittgenstein, that the $\sigma$ parameter
itself may be an illegitimate formalisation of something that resists
quantification.  We flag this as a foundational question for future work on the
epistemology of specification, rather than claiming to resolve it.

For AI coding agents that operate through conversation (iterative prompting,
evaluation, and revision), Suchman's critique is particularly sharp: the
``specification'' is not a document but an evolving dialogue whose informational
content is not fixed at any single moment.  The framework applies most naturally
when specifications are written documents with stable content, and less
naturally to interactive, conversational development where meaning is negotiated
in real time.  This limitation is genuine and not resolved by the current
framework.


\subsection{Governance Is Overhead}
\label{sec:contrarian-overhead}

\paragraph{The steelmanned argument.}
Governance has a cost curve that intersects its benefit curve at a specific
scale, and below that scale the cost exceeds the benefit.  For a two-person
startup, time spent writing ADRs and configuring fitness functions is time not
spent shipping features.  The strongest version: governance mechanisms are a
form of organisational scar tissue, growing in response to past failures but
imposing ongoing friction on teams that have already internalised the lessons.

Conway's Law~\citep{conway1968} strengthens this position: if architecture
mirrors organisational structure, and organisational metrics predict
failure-proneness with approximately 85\% precision~\citep{nagappan2008}
(outperforming all code metrics), then code-level governance treats a symptom.
Restructuring teams would be more effective than adding governance layers.  The
argument extends to agent teams: if agents mirror the organisational structures
of the teams that build them, then governing agent code is futile without
governing the organisational context that shapes agent behaviour.

\paragraph{What the evidence shows.}
Governance cost is roughly fixed per mechanism; benefit scales with team size
and codebase complexity.  The crossover appears at 5 to 10 developers, or more
precisely, at the point where no single person holds the complete mental model.
For distributed teams, formal governance becomes valuable earlier because
informal channels (hallway conversations, whiteboard sessions) are degraded or
absent.

Conway's Law admits an inverse: the inverse Conway
maneuver~\citep{skelton2019} deliberately restructures teams to produce desired
architecture.  Code-level governance and organisational governance are
complementary, not substitutes.  The novel question for AI agents is whether
agent governance requires organisational awareness, and the answer appears to be
yes: agent topology should mirror the desired system decomposition (the inverse
Conway maneuver applied to agents).

The governance capacity bound (Theorem~\ref{thm:capacity}) provides the formal
resolution: governance is overhead precisely when the drift rate is below the
bifurcation threshold.  Below the threshold, the system self-corrects without
governance.  Above it, governance is not overhead but a structural necessity;
its absence leads to coherence collapse.  The question is empirical: where does
a given team sit relative to the threshold?

The AI velocity corollary sharpens this analysis: when agents generate code at
10 to 100$\times$ the rate of human developers, the drift rate increases
proportionally while governance capacity remains constant (it is bounded by
human review bandwidth).  This pushes teams above the bifurcation threshold
earlier and more frequently.  The automated governance necessity follows:
governance must itself be partially automated to scale with agent throughput,
which is precisely what the multi-granularity cascade control architecture
provides.


\subsection{ADRs and Documentation Are Theater}
\label{sec:contrarian-adrs}

\paragraph{The steelmanned argument.}
ADRs can drift from purpose, enable responsibility diffusion, and suffer the staleness problem. Despite recommendations from ThoughtWorks, AWS, Microsoft, and Google, no published empirical study demonstrates that ADRs improve measurable software quality outcomes (defect rates, time-to-market, maintenance cost). The evidence is entirely practitioner reports and observational accounts.

The false confidence problem is acute: documentation with the appearance of completeness but no mechanism for detecting its own staleness creates a worse situation than acknowledged ignorance. Teams may defer to stale ADRs rather than exercising judgment, producing decisions that are formally ``compliant'' but substantively wrong.

\paragraph{What the evidence shows.}
ADRs provide clear value for onboarding and cross-team alignment. The primary failure mode is absent content (decisions made without ADRs) or bloated content (too many decisions documented), not wrong content. The survival analysis (Section~\ref{sec:survival}) provides a concrete remedy: tracking decision age alongside decision content, with domain-specific half-lives triggering proactive review.

Evidence expiry and the controlled relaxation operator $\delta$ address the staleness critique directly. Documentation with explicit uncertainty markers (``confidence: medium,'' ``last validated: 2025-03'') addresses the false confidence risk. The contribution is making these mechanisms formal rather than ad hoc.

The empirical gap is real and should be acknowledged honestly. The strongest evidence is that the \emph{problem} is real (AI accelerates architectural drift and quality degradation). The evidence that documentation and governance \emph{solve} it remains largely anecdotal. Designing governance systems to be testable, by tracking drift trajectories and decision decay, is a prerequisite for future empirical validation.

The ``ratchet is too rigid'' objection~\citep{ford2017} applies here as well: monotonically increasing documentation strictness prevents legitimate architectural evolution. The strongest governance systems already incorporate relaxation mechanisms. Fitness functions can be deprecated. Evidence expiry provides a natural clock for rule relaxation. The ratchet critique applies most forcefully to systems that treat rules as permanent rather than contextual, which is precisely what the lattice descent operator $\delta$ addresses through controlled relaxation with evidence expiry.


\subsection{The Human Should Just Review Everything}
\label{sec:contrarian-human}

\paragraph{The steelmanned argument.}
Human review is the gold standard for code quality. Rather than building elaborate multi-layer verification architectures, organisations should invest in thorough human review of all agent-generated code. Human reviewers can detect subtle architectural violations, naming convention drift, unnecessary abstractions, and strategic misalignment that no automated system can match. The four-layer defense architecture adds complexity and cost without matching the capabilities of a skilled human reviewer.

Mode selection is hard to automate~\citep{bainbridge1983ironies}: deciding when to apply lightweight versus heavyweight governance requires the kind of contextual judgment that resists formalisation. Bainbridge's ironies of automation create a genuine double bind: the more we automate governance decisions, the less capable humans become at making them when automation fails.

\paragraph{What the evidence shows.}
The throughput mismatch is the decisive counterargument. AI agents now generate over 50\% of committed code on major platforms. GitClear documents 211 million changed lines with an 8$\times$ increase in code duplication and a 60\% collapse in refactoring activity~\citep{gitclear2025}. Human review cannot scale to match this generation rate. At \$80 per PR for thorough human review (Layer~4 cost), the economics are prohibitive for high-throughput agent pipelines.

The detection ceiling theorem (Theorem~\ref{thm:ceiling}) shows that each layer has hard upper bounds on what it can detect, determined by available context. Human review has the highest ceiling but the lowest throughput. Layers~1 through 3 have lower ceilings but orders-of-magnitude higher throughput. The optimal architecture uses cheap, fast layers to filter the stream before expensive human attention is applied.

The Bayesian mode selection framework (Theorem~\ref{thm:mode}) provides a principled threshold for escalation: the asymmetric cost structure ($c_{\mathrm{FG}} \gg c_{\mathrm{GF}}$) means choosing the governed (human-reviewed) mode at even 5 to 10\% posterior probability of needing it is Bayes-optimal. This is not a replacement for human judgment but a triage mechanism that directs human attention where it has the highest marginal value.

The Dodge-Romig sampling plans~\citep{dodge1929method} and MIL-STD-105E switching rules provide the classical precedent: adaptive inspection regimes that tighten or relax based on observed quality. Our four-layer architecture adapts this paradigm to software verification, where the ``inspector'' at each layer has different capabilities and costs. The human reviewer remains indispensable for the undecidable slop types (meaningless tests, architectural erosion, boilerplate inflation) that resist automated detection, but the human's scarce attention is directed by the preceding automated layers rather than applied uniformly.


\subsection{AI Code Quality Is Good Enough}
\label{sec:contrarian-quality}

\paragraph{The steelmanned argument.}
The quality concerns are overstated. GitHub's controlled RCT~\citep{github2024rct} found that AI-assisted developers completed tasks 55\% faster with no measurable quality difference. For the vast majority of business software (CRUD applications, internal tools, prototypes), ``good enough'' quality is the economically rational target. The optimal quality level is explicitly $q^* < 1$ for low-stakes code; investing in elaborate quality architectures for code that will be rewritten in six months is waste.

Speed itself has value. The Yerkes-Dodson objection (``speed kills quality''~\citep{yerkes1908relation}) applies to human cognition under time pressure, but AI agents do not experience cognitive fatigue. The DORA finding that high-performing teams achieve both speed and stability~\citep{forsgren2018accelerate} suggests that the speed-quality tradeoff is not fundamental.

LLM-as-Judge is unreliable: inter-rater agreement (omega of 0.462 to 0.803) is concerning~\citep{llmjudgetrust2024}. Same-model judging amplifies correlated errors via the FKG inequality (Theorem~\ref{thm:fkg}). If the verification layer itself is unreliable, the entire quality architecture rests on shaky foundations.

\paragraph{What the evidence shows.}
The GitHub RCT measured short-term task completion, not long-term maintainability. GitClear's longitudinal data tells the opposite story: 17\% maintainability drops, refactoring declining from 25\% to less than 10\% of changed lines, and bug rates rising 12\% at six months~\citep{gitclear2025}. The compound degradation theorem (Theorem~\ref{thm:degradation}) formalises why: individually CI-passing changes can produce superlinear entropy growth through accretion effects that no single-change review can detect.

The ``slop is acceptable'' position is partially validated by our own framework. The optimal quality level (Theorem~\ref{thm:optquality}) is explicitly $q^* < 1$ for low-stakes code. The contribution is making this precise: for prototypes, Layer~1 only; for production systems, all four layers. The framework does not argue for maximal quality everywhere; it argues for type-specific, cost-aware quality investment.

On LLM-as-Judge reliability: same-model judging is indeed problematic (Theorem~\ref{thm:fkg}). But cross-model judging with structural scaffolding (Layer~3 constrained by Layers~1 and 2) remains cost-effective for semi-decidable types. The diversity gain (Definition~\ref{def:diversity}) shows that a weak independent detector outperforms a strong correlated one, providing a concrete architectural remedy.

VIH (Verified Iterations per Hour) captures the speed-quality tradeoff more precisely than raw throughput: $\VIH = \lambda_{\mathrm{raw}} \cdot p_{\mathrm{verify}}$. The strongest form of the ``speed kills quality'' objection is that verification itself becomes less reliable under time pressure, creating a feedback loop that systematically overestimates VIH at high speeds. This remains an open empirical question that the framework does not resolve.

VIH is also not a complete objective function. A verified iteration that introduces subtle architectural debt, passes tests but degrades maintainability, or solves the wrong subproblem is ``verified'' but not valuable. Reinertsen's cost-of-delay framework provides a more comprehensive alternative that accounts for the economic value of delivery timing. VIH is a useful proxy for the speed-quality tradeoff specifically, but it should not be mistaken for a measure of total value delivered. The ``18 types are not the right taxonomy'' objection from the quality pillar applies here as well: the three-factor latent structure (context blindness, completion bias, training leakage) may be more fundamental than the 18-type enumeration; confirmatory factor analysis on appropriate datasets would test this.

The ``caching is a false economy'' objection deserves particular attention. Cache invalidation bugs are intermittent, state-dependent, and compound through cascading cached values. For rapidly changing codebases under active development, cache hit rates may be too low to justify the infrastructure overhead. Our EOQ theorem provides a principled refresh schedule, but assumes knowledge of the change rate $\lambda$, which itself must be estimated. The ``fresh eyes'' advantage (17.1\% quality improvement from fresh context) is large enough that naive persistent state is VIH-negative; incremental context is only justified when compaction reduces quality loss below 12.3\%.


\subsection{Naur Is Right and Theory Cannot Be Preserved}
\label{sec:contrarian-naur}

\paragraph{The steelmanned argument.}
Naur~\citep{naur1985} argued that programming is theory building, where ``theory'' (in Ryle's sense) is the knowledge a person must have ``in order not only to do certain things intelligently but also to explain them, to answer queries about them, to argue about them.'' Three capabilities resist documentation: mapping between code and reality, design justification (``the justification is and must always remain the programmer's direct, intuitive knowledge or estimate''), and modification intelligence, which depends on recognising similarities ``that are not, and cannot be, expressed in terms of criteria.''

The tacit divergence conjecture (Conjecture~\ref{conj:tacit}) formalises this: for Kolmogorov-random components of tacit knowledge, the description length diverges, meaning that no finite documentation can capture them. If this is correct, then THEORY.md, ADRs, and all documentation artifacts are lossy compression creating false confidence. Acting on extracted ``decisions'' may be worse than no documentation at all, because it creates the illusion of understanding where none exists.

Tsoukas's critique of Nonaka and Takeuchi deepens the problem: tacit and explicit knowledge are not convertible pools but two dimensions of all knowing. The ``conversion'' metaphor (tacit to explicit and back) is incoherent; what actually happens is articulation, which necessarily transforms and reduces the original knowledge.

\paragraph{What the evidence shows.}
Naur wrote in 1985 about small teams with long tenure. Modern software involves high turnover across time zones and continents. The alternative to lossy documentation is not ``direct contact with theory-holders'' but \emph{no knowledge transfer at all}. Lossy compression is still enormously useful: JPEG images demonstrate this. The question is not whether documentation is perfect but whether it is better than nothing, and whether its limitations are understood.

The double bottleneck theorem (Theorem~\ref{thm:bottleneck}) provides the formal resolution: governance fidelity is bounded by log richness, and agent logs that externalise decisions, assumptions, and rejected alternatives capture more theory than code alone, even if the capture is necessarily lossy. Documentation with explicit uncertainty markers (``confidence: medium,'' ``last validated: 2025-03'') addresses the false confidence risk by making the lossiness visible rather than hiding it.

The survival analysis (\S\ref{sec:survival}) adds a temporal dimension: decision validity decays with domain-specific half-lives (UI/frontend decisions, approximately 1.2 years; API contracts, approximately 2.5 years; data model decisions, approximately 4.3 years). Governance that tracks decision age and proactively surfaces decisions approaching their expected half-life provides a practical response to Naur's concern, acknowledging that documentation decays while building mechanisms to detect and respond to that decay.

We accept the residual: perfect externalisation is impossible. The framework uses this impossibility as a design constraint (Design Principle~11: accept residual theory loss), supplementing artifact-based governance with social mechanisms (code review conversations, architecture advice processes) for knowledge that resists externalisation.

The Dreyfus model of skill acquisition provides additional nuance: expert
practitioners operate on intuitions that resist rule-based decomposition.
Collins's taxonomy of tacit knowledge (relational, somatic, collective)
identifies which components are in principle externalisable (relational) and
which are not (somatic, collective).  The governance architecture should invest
extraction effort proportional to externalisability, rather than attempting
uniform documentation.  This selective approach respects Naur's insight while
remaining practically useful.


\subsection{Cross-Cutting Assessment}

Several patterns emerge across these seven positions.  First, most contrarian
arguments are \emph{partially} correct rather than wholly right or wrong.  The
framework's response is typically to identify the conditions under which the
objection holds ($\sigma < 0.3$ for the formalism critique, below 5 to 10
developers for the governance-as-overhead critique, $q^* < 1$ for the
quality-is-good-enough critique) and to incorporate those conditions as design
parameters rather than dismissing the objection.

Second, the strongest objections challenge the framework's \emph{foundations}
(Suchman on situated action, Wittgenstein on language games, Naur on
theory preservation) rather than its \emph{conclusions}.  We acknowledge these
as genuine limitations that define the framework's scope of applicability rather
than falsify its results within that scope.

Third, the empirical evidence is strongest for the \emph{problems} the framework
addresses (compound errors in pipelines, quality degradation from AI code
generation, governance scaling challenges) and weakest for the \emph{solutions}
it proposes (ADRs improving quality outcomes, formal methods reducing
maintenance costs, governance preventing drift).  This asymmetry is honest but
should constrain confidence in the prescriptive claims.


% ============================================================
% RELATED WORK (insert before Conclusion)
% ============================================================
\section{Related Work}
\label{sec:related}

We position this work against six bodies of prior research.  Each subsection
identifies the specific contributions we build upon and clarifies where our
framework extends, synthesises, or merely applies existing results.


\subsection{Formal Methods in Software Engineering}

The refinement lattice (Section~\ref{sec:refinement}) builds directly on Back's
refinement calculus~\citep{back1980correctness} and Morgan's \emph{Programming
from Specifications}~\citep{morgan1994programming}, which formalised the
stepwise refinement of specifications into code.  Hoare and He's Unifying
Theories of Programming~\citep{hoare1998unifying} showed that programs are
predicates and refinement is logical implication, placing specifications and
code in a single semantic universe.  Abrial's B-Method and
Event-B~\citep{abrial2010modeling} provide the most industrially deployed
refinement-based formal method, with verified systems including the Paris
M\'{e}t\'{e}or driverless line, directly instantiating the tower of Galois
connections.

Cousot and Cousot's abstract interpretation
framework~\citep{cousot1977abstract} formalised abstraction via Galois
connections five decades ago, leading to mature industrial tools (Astr\'{e}e,
Polyspace, Infer).  We use the Galois connection formalism to model the
abstraction spectrum but do not extend the abstract interpretation framework
itself.  Roy and Cordy's clone type taxonomy (Types~I through~IV) characterises
detection decidability for duplicate logic: Types~I through~III (textual,
renamed, gapped) are decidable; Type~IV (semantic clones) is not.  This
distinction directly informs our decidability classification of slop types.

Murphy, Notkin, and Sullivan's reflexion models~\citep{murphy1995, murphy2001}
compare intended architecture against extracted source models, providing the
formal basis for architectural erosion detection and the conformance checking
that underpins our governance pillar.  The Curry-Howard-Lambek
correspondence~\citep{wadler2015propositions} provides the type-theoretic
foundation connecting proofs, programmes, and categories that unifies our
treatment of specifications as both logical propositions and programme types.

For AI-assisted formal methods, Kleppmann~\citep{kleppmann2025formal} predicted
that AI would make formal verification mainstream;
Congdon~\citep{congdon2025formal} endorsed and extended this argument with a
concrete workflow vision.  The vericoding benchmark~\citep{vericoding2025}
provides early empirical support.  Our framework formalises \emph{why} this
convergence is expected: AI reduces the cost of traversing the refinement
lattice, while formal verification provides a deterministic oracle.


\subsection{Multi-Agent Systems and Coordination}

Kim et al.~\citep{kim2025scaling} provide the most comprehensive empirical
evaluation of multi-agent system scaling, establishing error amplification
factors, capability saturation thresholds, and communication scaling laws across
five canonical architectures.  Our work builds on their empirical findings by
developing a formal framework that explains their observations and derives
actionable design rules.  Cemri et al.~\citep{Cemri2025} introduced MAST, the
first empirically grounded failure taxonomy for multi-agent LLM systems,
identifying 14 failure modes across over 1,600 traces.

For merge conflict research, Brun et al.~\citep{Brun2011} pioneered proactive
conflict detection with Crystal.  Kasi and Sarma~\citep{KasiSarma2013}
developed Cassandra for conflict-minimising task scheduling.  Ghiotto
et al.~\citep{Ghiotto2018} provided the largest study of merge conflict
characteristics (2,731 Java projects).  Sousa et al.~\citep{Sousa2018}
formalised semantic conflict-freedom verification with SafeMerge.  Zhu
et al.~\citep{Zhu2024congra} introduced ConGra, finding (counterintuitively)
that general-purpose LLMs outperform code-specialised ones at conflict
resolution, a result relevant to coordination agent design.  For merge tools,
Falleri et al.~\citep{Falleri2014} developed GumTree for fine-grained AST
differencing; Shen et al.~\citep{Shen2019} developed IntelliMerge for
refactoring-aware merging (88.48\% precision); and Apel et al.~\citep{Apel2012}
developed JDime for semistructured merge (39\% conflict reduction).

Graph partitioning, which formalises our task decomposition problem, has deep
roots: Fiedler~\citep{Fiedler1973} introduced algebraic connectivity and
spectral bisection; Kernighan and Lin~\citep{KernighanLin1970} developed the KL
heuristic; Fiduccia and Mattheyses~\citep{FiducciaMattheyses1982} improved it
with linear-time passes and hypergraph support; Karypis and
Kumar~\citep{KarypisKumar1998} developed METIS with $O(|E|)$ multilevel
partitioning; Arora, Rao, and Vazirani~\citep{AroraRaoVazirani2009} achieved
$O(\sqrt{\log n})$ sparsest cut via SDP; and Lee et al.~\citep{Lee2012cheeger}
proved higher-order Cheeger inequalities connecting spectral gaps to partition
quality.

Contract-based design provides our assume-guarantee framework:
Meyer~\citep{Meyer1992} introduced Design by Contract; Henzinger
et al.~\citep{Henzinger2002} developed assume-guarantee methodology;
Jones~\citep{Jones1983} developed rely-guarantee reasoning; Benveniste
et al.~\citep{Benveniste2018} provided a comprehensive theory of contracts for
system design.  Our coordination contracts extend these to the specific setting
of LLM agents operating on shared codebases.

For distributed systems foundations, Bernstein and Hadzilacos~\citep{BernsteinHadzilacos1987} established concurrency control theory; Berenson et al.~\citep{Berenson1995} formalised isolation levels and identified write skew; Fekete et al.~\citep{Fekete2005} showed how to make snapshot isolation serialisable via promotion, analogous to our use of contracts to close the gap from snapshot to serialisable isolation for agents. Lamport~\citep{Lamport1978} defined happened-before and logical clocks; Herlihy~\citep{Herlihy1991} proved the consensus number hierarchy; and the FLP impossibility result~\citep{FLP1985} proved that deterministic consensus is impossible with one faulty process. Ongaro and Ousterhout~\citep{Ongaro2014} developed Raft for understandable consensus. These results inform our treatment of coordination correctness: the strong correctness condition (serialisable merge outcomes) is analogous to linearisability, while the weak correctness condition (compiles and passes tests) provides a practically achievable alternative.

For software modularity, Mitchell and Mancoridis~\citep{Mitchell2006bunch} developed Bunch for automatic software modularisation. Martin~\citep{Martin1994} defined coupling metrics (afferent, efferent, instability) that underpin our coupling graph formalism. MacCormack et al.~\citep{MacCormack2012} validated Conway's Law empirically via the mirroring hypothesis, confirming that organisational structure predicts architectural structure with high fidelity.


\subsection{Software Quality and Testing}

Software defect taxonomies provide context for our slop type classification. IEEE~1044~\citep{ieee1044} provides multi-attribute anomaly classification. ODC~\citep{chillarege1992} offers eight orthogonal defect types with development-phase correlations. Beizer~\citep{beizer1990} provides a ten-category bug taxonomy. CWE covers security vulnerabilities. None of these captures accretion defects (Section~\ref{sec:degradation}), the novel defect class that is correct, unnecessary, individually defensible, and collectively harmful.

For software reliability, Musa's models~\citep{musa1987} connect operational profiles to defect discovery. Fenton and Neil~\citep{fenton1999} introduced Bayesian networks for defect prediction. Capers Jones~\citep{capersjones2012} provided the foundational defect removal effectiveness data from 25,000 projects. Littlewood and Strigini~\citep{littlewood1993, littlewood2000} proved that diverse software versions do not fail independently, a result we adapt to LLM judge-producer pairs via the FKG inequality.

The classical reliability theory foundations include Barlow and Proschan's series system product law and importance measures~\citep{barlow1965mathematical, barlow1975statistical}, the IFR/IFRA closure theorems~\citep{birnbaum1961stochastic}, and the Young-Daly formula~\citep{young1974first, daly2006higher} for optimal checkpoint intervals. The Jelinski-Moranda~\citep{jelinski1972software} and Musa-Okumoto~\citep{musa1984logarithmic} models of fault detection as a depletion process inform our verification convergence analysis.

For AI code quality measurement, CodeRabbit~\citep{coderabbit2025}, GitClear~\citep{gitclear2025}, Qodo~\citep{qodo2025}, and the ``debt boom'' analysis~\citep{debtboom2026} provide the emerging empirical base. The controlled GitHub trial~\citep{github2024rct} stands in tension with observational results, a contradiction we interpret through the verification bottleneck (short-term task completion gains masking long-term maintainability costs). The Swiss Cheese Model~\citep{reason1990} frames accident causation through aligned holes in multiple barriers. Viola-Jones cascades~\citep{violajones2004} formalise sequential classifier design with early rejection of easy negatives. Boehm's cost escalation curve~\citep{boehm1981} and Capers Jones~\citep{capersjones2012} provide the economic foundation for layered defense, though the cost curve may be flatter in modern CI/CD environments.

Anthropic's agent evaluation methodology~\citep{anthropic2026evals} introduces the pass@$k$ vs.\ pass$^k$ distinction, clarifying whether verification requires any correct output in $k$ attempts or correct outputs across all $k$ attempts. The LLM-as-Judge survey~\citep{li2024survey} documents systematic biases (sycophancy, position bias) that limit LLM-based verification. The METR merge-gap study~\citep{metr2026swebench} provides the most striking evidence of circular validation bias: agents that ``solve'' benchmarks by passing tests frequently produce code that human reviewers would reject.


\subsection{Architecture Governance and ADRs}

Nygard~\citep{nygard2011} proposed the lightweight ADR format. The Architecture Advice Process~\citep{harmellaw2023} adds a governance layer: anyone can make architectural decisions, but must first consult affected parties and domain experts. Ford et al.~\citep{ford2017} defined architectural fitness functions as objective integrity assessments of architectural characteristics, with tools like ArchUnit encoding these as executable tests. The gap in ADR-to-fitness-function automation is being addressed by tools like Archgate through the ratchet model.

For reflexion models, Murphy et al.~\citep{murphy1995, murphy2001} provide the formal foundation for conformance checking, with extensions including behavioural reflexion models and CI-integrated prototypes detecting nonconformances ``within minutes, instead of the months or years it takes today.'' The AgDR format extends ADRs with agent-specific metadata (agent name, model ID, session, trigger type), addressing traceability gaps when AI agents make architectural choices.

Conway's Law~\citep{conway1968} established the homomorphism between organisational and system structure. Nagappan et al.~\citep{nagappan2008} found that organisational metrics predict failure-proneness with approximately 85\% precision and recall, outperforming all code metrics. MacCormack et al.~\citep{maccormack2008} validated the mirroring hypothesis empirically. Skelton and Pais~\citep{skelton2019} operationalised the inverse Conway maneuver through Team Topologies.

Hellerstein et al.~\citep{hellerstein2004} provided the foundational treatment of feedback control for computing systems. Our contribution is applying cascade control specifically to multi-granularity governance, with stability analysis and Nyquist-like sampling requirements for governance loops. Lehman's laws of software evolution~\citep{Lehman1980}, particularly the law of increasing complexity, underpin our coupling density dynamics model.

The Goodhart's Law literature~\citep{goodhart1975, manheim2019goodhart} is directly relevant to governance measurement. When a governance metric becomes a target, it ceases to be a good measure: coupling targets incentivise restructuring that reduces measured coupling while introducing unmeasured coupling through shared configuration; coverage targets incentivise trivial fitness functions. The defense is metric rotation, metric triangulation, and cultural alignment. Our framework addresses this through the multi-granularity approach: different governance loops measure different quantities at different frequencies, making systematic gaming more difficult.

The DORA programme~\citep{dora2025} provides the largest objective measurement of AI's impact on software engineering practice. The productivity paradox it documents (individual gains, organisational stagnation) motivates our quality-adjusted speedup metric, which formalises the quality cost of parallelism. McKinney~\citep{McKinney2026} argued that Brooks's Law applies directly to LLM agents, a position our Extended Amdahl's Law formalises.


\subsection{Information Theory Applied to Software}

The application of algorithmic information theory to software has precedent. Li and Vit\'{a}nyi~\citep{li2019introduction} discuss the connection between algorithmic complexity and programme equivalence. The normalised information distance~\citep{li2004similarity} has been applied to software clone detection and plagiarism, though not (to our knowledge) to the specification-code relationship specifically. Our abstraction gap $\mathcal{G}(S, P) = K(P \mid S)$ applies conditional Kolmogorov complexity to this relationship.

Hindle et al.~\citep{hindle2012naturalness} established that code is ``natural'' (statistically predictable), measuring code entropy and showing that language models can exploit this regularity. Hellendoorn and Devanbu refined these measurements, finding entropy of approximately 1.25 bits per token. Shannon's original measurements~\citep{shannon1948} establish the baseline at 0.6 to 1.3 bits per character for natural language. These measurements ground our information-theoretic treatment of the specification-code channel.

Hassan's change entropy~\citep{hassan2009} predicts faults from change scattering. Our codebase entropy (Definition~\ref{def:entropy}) shifts from measuring change scattering to measuring the coupling structure that necessitates it. Tishby et al.'s information bottleneck method~\citep{tishby2000information} provides an alternative formalisation: the context window as an information bottleneck between the full codebase and the generated code.

For context engineering specifically, Anthropic~\citep{anthropic2025context} established the ``smallest set of high-signal tokens'' principle with four functional components (system prompts, tools, examples, message history). Spotify's Honk~\citep{spotify2025honk} contributed the ``static over dynamic'' philosophy. HumanLayer~\citep{humanlayer2025skillissue} distinguished structural from instructional enforcement and coined ``success should be silent.'' Lin and Bilmes~\citep{lin2011submodular} proved submodularity for document summarisation functions. Jina et al.~\citep{jina2025submodular} independently applied submodular optimisation to LLM context engineering, validating our formalisation. Liu et al.~\citep{liu2024lost} identified the ``lost in the middle'' phenomenon. Subsequent work by Chroma~\citep{chroma2025contextrot}, Du et al.~\citep{du2025context}, Kuratov et al.~\citep{kuratov2024babilong}, and Hsieh et al.~\citep{hsieh2024ruler} collectively establishes the empirical degradation landscape that our context degradation function formalises.

The value of information literature, originating with Howard (1966), frames context selection as a decision-theoretic problem: which information to acquire before acting. The greedy algorithm for submodular value of information was analysed by Krause and Guestrin (2005). Our degradation penalty is a novel addition absent from the classical setting, where acquiring more information is assumed to be never harmful. For code, this assumption fails: code dependencies create non-monotonic relevance effects where additional context can actively mislead.


\subsection{AI Coding Agents and Harnesses}

The empirical landscape of AI coding agents provides both motivation and calibration for our framework. SWE-bench~\citep{sweeffi2025} measures resolve rates, with SWE-Effi introducing resource-bounded effectiveness scores. Bian et al.~\citep{bian2025limits} decompose bottlenecks, finding that web environment latency dominates (53.7\% of runtime) with LLM API latency varying up to 69$\times$. These findings motivate our temporal architecture.

For speculative execution in agent pipelines, Zhou et al.~\citep{zhou2025speculative} achieve 46 to 55\% next-action prediction accuracy; Ji et al.~\citep{ji2026dualspec} extend this with heterogeneous speculation, achieving 3.28$\times$ speedup. Our contribution is the formal speculation ratio test and depth limit, grounded in classical scheduling theory and connected to the Spectre analogy for hidden costs.

Agent verification frameworks include AgentSpec~\citep{wang2026agentspec}, which provides the key empirical data point on structural versus prompt enforcement (87\% versus 77\%); AgentGuard~\citep{koohestani2025agentguard}, which introduces dynamic probabilistic assurance via MDP modelling; and Multi-Agent Verification~\citep{lifshitz2025mav}, which demonstrates that diverse verifier ensembles outperform homogeneous ones and that weaker models can verify stronger ones.

For scalable oversight, Irving et al.'s debate framework~\citep{irving2018debate} proposes using adversarial argumentation between agents to help weaker judges evaluate stronger systems. Constitutional AI~\citep{bai2022constitutional} uses self-critique against a set of principles; while effective as a training methodology, it inherits the circular validation limitations we formalise. METR's reward-hacking analysis~\citep{metr2025rewardhacking}, Apollo Research's scheming studies~\citep{apollo2025scheming}, and NIST CAISI's cheating documentation~\citep{nist2025cheating} collectively establish that structural enforcement is necessary for adversarial robustness.

For caching and context persistence, Liang et al.~\citep{liang2026cache} evaluate prompt caching empirically, and Zhang et al.~\citep{zhang2025apc} introduce plan template caching at the reasoning level. Suzgun and Kalai~\citep{suzgun2024metaprompting} establish the ``fresh eyes'' advantage that motivates our fresh context dominance theorem. The speed-quality tradeoff literature includes Reinertsen's cost-of-delay framework, Forsgren et al.'s~\citep{forsgren2018accelerate} evidence that speed and stability can improve simultaneously through DevOps practices, and Kahneman's~\citep{kahneman2011thinking} cognitive science foundation for why speed pressure degrades complex reasoning.

Industrial harness implementations provide the strongest empirical grounding. Cursor~\citep{Cursor2026} documented the evolution from flat locking to planner/worker/judge topology, providing evidence that hierarchical role separation succeeds at scale. Osmani~\citep{Osmani2025} synthesised practitioner patterns including optimal team sizing and tier-based tool stacks. ChatDev, MetaGPT, and AgentCoder explore role-based agent decomposition for software tasks, providing additional data points on multi-agent coordination architectures. The convergence of these production systems toward architectures predicted by our framework (layered verification, structural enforcement, bounded context windows, governance loops) provides indirect validation, though controlled experiments remain necessary.

The sampling inspection literature from manufacturing quality control provides precedent for our adaptive verification approach. The Dodge-Romig sampling plans~\citep{dodge1929method, dodge1959sampling} and MIL-STD-105E switching rules~\citep{milstd105e} formalised adaptive inspection regimes that tighten or relax based on observed quality. The Lindsay-Bishop dynamic programming formulation for inspection allocation~\citep{lindsay1964allocation} provides the discrete optimisation framework that informs our verification scheduling result. These classical results, developed for physical manufacturing, transfer naturally to software verification pipelines where the ``inspection'' is a verification step and the ``product'' is generated code.

\paragraph{Position in the landscape.}
This paper is a synthesis and formalisation, not a claim of novel mathematical results. The information-theoretic machinery is standard~\citep{li2019introduction}; the refinement calculus is established~\citep{back1980correctness, morgan1994programming}; the reliability theory is classical~\citep{barlow1965mathematical}; the Curry-Howard correspondence is a theorem~\citep{wadler2015propositions}. Our contribution is the application of these tools to the specific problem of AI coding agent harness design, the unification across seven architectural pillars through shared formal machinery, and the identification of cross-pillar constraints (the abstraction gap chain, the compound error cascade, the governance capacity bound) that emerge only from the unified treatment.
